\documentclass[11pt]{article}
\input{../../shared/project_style.tex}
\rhead{Basics 36 Textbook Draft}

\begin{document}
\DocTitle{Basics 36: Coupled Harmonics and Avoided Crossings}{V - Plasma waves and Alfvén physics}{Two-mode matrices, gap formation, harmonic mixing, positivity, and mode tracking}

\begin{overviewbox}
\textbf{Textbook draft, version 1.0.} Toroidal and stellarator geometry couples Fourier harmonics that would be independent in a uniform medium. The simplest two-harmonic model already explains avoided crossings, spectral gaps, mixed eigenvectors, and the numerical issues central to the project's reduced Module 4 benchmark. This chapter is designed for a reader with no prior plasma-physics training. It distinguishes exact identities from physical approximations and connects the central equations to later simulation modules.
\end{overviewbox}

\ProjectTOC

\section{Learning objectives}
After completing this note, the reader should be able to:
\begin{itemize}[leftmargin=1.6em]
\item Solve a two-mode coupled eigenvalue problem analytically.
\item Distinguish exact crossing from avoided crossing.
\item Calculate gap width and harmonic mixing.
\item Explain positivity constraints on a squared-frequency operator.
\item Connect local coupling to a global radial eigenproblem.
\item Track modes robustly across parameter scans.
\end{itemize}

\section{Prerequisites and reading strategy}
\begin{itemize}[leftmargin=1.6em]
\item Basics 6, 7, 34, and 35.
\end{itemize}
On a first reading, emphasize physical meaning and dimensional checks. On a second reading, reproduce the derivations. Attempt the computational laboratory only after the limiting cases are understood.

\section{Uncoupled two-mode model}

Let two local squared frequencies be $a(s)=\omega_1^2(s)$ and $b(s)=\omega_2^2(s)$. Without coupling,
\begin{equation}
\mathsf H_0=
\begin{bmatrix}
a&0\\
0&b
\end{bmatrix}.
\end{equation}
The eigenvalues cross wherever $a=b$. The corresponding eigenvectors remain pure basis harmonics.


\section{Coupled matrix and eigenvalues}

Introduce symmetric coupling $g(s)$:
\begin{equation}
\mathsf H=
\begin{bmatrix}
a&g\\
g&b
\end{bmatrix}.
\end{equation}
Its squared-frequency eigenvalues are
\begin{equation}
\boxed{\lambda_\pm=
\frac{a+b}{2}\pm
\sqrt{\left(\frac{a-b}{2}\right)^2+g^2}}.
\end{equation}
For $g\ne0$, $\lambda_+\ne\lambda_-$ even at $a=b$. The exact crossing is replaced by an avoided crossing.


\section{Gap at the uncoupled crossing}

At $a=b=a_c$,
\begin{equation}
\lambda_\pm=a_c\pm|g_c|.
\end{equation}
The gap in squared frequency is $2|g_c|$. In ordinary frequency,
\begin{equation}
\Delta\omega=
\sqrt{a_c+|g_c|}-\sqrt{a_c-|g_c|},
\end{equation}
provided $a_c>|g_c|$. For weak coupling,
\begin{equation}
\Delta\omega\simeq\frac{|g_c|}{\sqrt{a_c}}.
\end{equation}
One must distinguish a gap in $\omega^2$ from a gap in $\omega$.


\section{Harmonic mixing and mixing angle}

Write normalized eigenvectors as
\begin{equation}
\mathbf v_+=
\begin{bmatrix}\cos\theta\\ \sin\theta\end{bmatrix},
\qquad
\mathbf v_-=
\begin{bmatrix}-\sin\theta\\ \cos\theta\end{bmatrix}.
\end{equation}
The angle satisfies
\begin{equation}
\tan(2\theta)=\frac{2g}{a-b}.
\end{equation}
Far from crossing, each mode resembles one uncoupled harmonic. Near crossing, $\theta\simeq\pi/4$ and both harmonics contribute comparably. Following modes only by ordered eigenvalue therefore makes their apparent identity exchange across the avoided crossing.


\section{Positivity-preserving coupling}

If $\mathsf H$ represents squared frequencies, it should be positive semidefinite in the intended stable local benchmark. For a symmetric $2\times2$ matrix,
\begin{equation}
a\ge0,\qquad b\ge0,\qquad g^2\le ab.
\end{equation}
A useful parameterization is
\begin{equation}
\boxed{g=\epsilon_c\sqrt{ab}},\qquad |\epsilon_c|<1.
\end{equation}
This automatically respects $g^2<ab$ where $a,b>0$ and causes coupling to vanish if either uncoupled branch reaches zero. The current repository uses this form to avoid an unphysical negative local squared frequency.


\section{From local matrix to radial differential operator}

A global mode has radial structure. One reduced form is
\begin{equation}
-\frac{d}{ds}\left(\mathsf K(s)\frac{d\boldsymbol\xi}{ds}\right)
+\mathsf H(s)\boldsymbol\xi
=\omega^2\mathsf M(s)\boldsymbol\xi.
\end{equation}
The diagonal potential terms reconstruct local continuum branches; off-diagonal terms couple harmonics; $\mathsf K$ produces radial communication and boundary quantization. A discrete eigenvalue inside a local gap can represent a global Alfvén eigenmode.


\section{Mode tracking across scans}

Eigenvalues sorted by magnitude can swap character near an avoided crossing. Track a mode using mass-weighted overlap,
\begin{equation}
\mathcal O_{ij}=
\frac{|\mathbf x_i^\dagger\mathsf M\mathbf y_j|}
{\sqrt{(\mathbf x_i^\dagger\mathsf M\mathbf x_i)
(\mathbf y_j^\dagger\mathsf M\mathbf y_j)}}.
\end{equation}
Maximum overlap, harmonic content, and radial localization together provide more reliable identity than eigenvalue index alone.


\section{Physical meaning in toroidal plasmas}

Toroidicity, ellipticity, and helical geometry introduce equilibrium Fourier coefficients that multiply perturbation harmonics and generate off-diagonal coupling. The two-mode model is not a complete stellarator spectrum, but it captures the universal mechanism: symmetry-breaking geometry couples branches, opens gaps, and permits discrete mixed-harmonic modes.


\section{Computational laboratory}
Scan a two-by-two matrix through an uncoupled crossing, plot eigenvalues and harmonic weights, and verify the analytic formulas. Compare constant coupling with positivity-preserving coupling. Then solve the repository's global two-harmonic radial benchmark and track modes under changes in coupling strength.

\section{Summary}
\begin{itemize}[leftmargin=1.6em]
\item Geometry-induced coupling converts exact branch crossings into avoided crossings.
\item The coupled eigenvalues and harmonic mixing follow analytically from a two-mode matrix.
\item Squared-frequency models require positivity-aware coupling choices.
\item Radial differential terms convert local gaps into global eigenmode problems.
\item Overlap and harmonic content are essential for tracking modes across parameter changes.
\end{itemize}

\SymbolsAcronymsSection
\begin{symtable}
$a,b$ & uncoupled squared frequencies & Diagonal two-mode entries. \\
$g$ & coupling coefficient & Off-diagonal harmonic interaction. \\
$\lambda_\pm$ & coupled squared frequencies & Eigenvalues of the local matrix. \\
$\theta$ & mixing angle & Harmonic composition parameter. \\
$\epsilon_c$ & dimensionless coupling strength & In $g=\epsilon_c\sqrt{ab}$. \\
$\mathcal O_{ij}$ & mode overlap & Mass-weighted eigenvector similarity. \\
\end{symtable}

\section{Exercises}
\begin{enumerate}[leftmargin=1.8em]
\item Derive the two coupled eigenvalues from the characteristic polynomial.
\item Derive the mixing-angle formula.
\item Convert a gap in squared frequency to a gap in frequency.
\item Prove the positive-semidefinite condition $g^2\le ab$.
\item Explain why eigenvalue sorting alone fails as a mode-tracking method.
\end{enumerate}

\section{References}
\begin{thebibliography}{99}
\bibitem{ref1} G. B. Arfken, H. J. Weber, and F. E. Harris, \emph{Mathematical Methods for Physicists}, 7th ed., Academic Press (2013).
\bibitem{ref2} G. Strang, \emph{Linear Algebra and Its Applications}, 4th ed., Brooks/Cole (2006).
\bibitem{ref3} W. A. Strauss, \emph{Partial Differential Equations: An Introduction}, 2nd ed., Wiley (2007).
\bibitem{ref4} J. P. Freidberg, \emph{Ideal MHD}, Cambridge University Press (2014).
\bibitem{ref5} J. P. Goedbloed, R. Keppens, and S. Poedts, \emph{Advanced Magnetohydrodynamics}, Cambridge University Press (2010).
\bibitem{ref6} H. Goedbloed and S. Poedts, \emph{Principles of Magnetohydrodynamics}, Cambridge University Press (2004).
\bibitem{ref7} T. H. Stix, \emph{Waves in Plasmas}, American Institute of Physics (1992).
\bibitem{ref8} D. A. Gurnett and A. Bhattacharjee, \emph{Introduction to Plasma Physics}, 2nd ed., Cambridge University Press (2017).
\bibitem{ref9} F. F. Chen, \emph{Introduction to Plasma Physics and Controlled Fusion}, 3rd ed., Springer (2016).
\bibitem{ref10} R. J. Goldston and P. H. Rutherford, \emph{Introduction to Plasma Physics}, CRC Press (1995).
\bibitem{ref11} P. M. Bellan, \emph{Fundamentals of Plasma Physics}, Cambridge University Press (2006).
\end{thebibliography}

\end{document}
